\documentclass[handout,nooutcomes,noauthor]{ximera}

\graphicspath{  
{./}
{./whoAreYou/}
{./drawingWithTheTurtle/}
{./bisectionMethod/}
{./circles/}
{./anglesAndRightTriangles/}
{./lawOfSines/}
{./lawOfCosines/}
{./plotter/}
{./staircases/}
{./pitch/}
{./qualityControl/}
{./symmetry/}
{./nGonBlock/}
}


%% page layout
\usepackage[cm,headings]{fullpage}
\raggedright
\setlength\headheight{13.6pt}


%% fonts
\usepackage{euler}

\usepackage{FiraMono}
\renewcommand\familydefault{\ttdefault} 
\usepackage[defaultmathsizes]{mathastext}
\usepackage[htt]{hyphenat}

\usepackage[T1]{fontenc}
\usepackage[scaled=1]{FiraSans}

%\usepackage{wedn}
\usepackage{pbsi} %% Answer font


\usepackage{cancel} %% strike through in pitch/pitch.tex


%% \usepackage{ulem} %% 
%% \renewcommand{\ULthickness}{2pt}% changes underline thickness

\tikzset{>=stealth}

\usepackage{adjustbox}

\setcounter{titlenumber}{-1}

%% journal style
\makeatletter
\newcommand\journalstyle{%
  \def\activitystyle{activity-chapter}
  \def\maketitle{%
    \addtocounter{titlenumber}{1}%
                {\flushleft\small\sffamily\bfseries\@pretitle\par\vspace{-1.5em}}%
                {\flushleft\LARGE\sffamily\bfseries\thetitlenumber\hspace{1em}\@title \par }%
                {\vskip .6em\noindent\textit\theabstract\setcounter{question}{0}\setcounter{sectiontitlenumber}{0}}%
                    \par\vspace{2em}
                    \phantomsection\addcontentsline{toc}{section}{\thetitlenumber\hspace{1em}\textbf{\@title}}%
                     }}
\makeatother



%% thm like environments
\let\question\relax
\let\endquestion\relax

\newtheoremstyle{QuestionStyle}{\topsep}{\topsep}%%% space between body and thm
		{}                      %%% Thm body font
		{}                              %%% Indent amount (empty = no indent)
		{\bfseries}            %%% Thm head font
		{)}                              %%% Punctuation after thm head
		{ }                           %%% Space after thm head
		{\thmnumber{#2}\thmnote{ \bfseries(#3)}}%%% Thm head spec
\theoremstyle{QuestionStyle}
\newtheorem{question}{}



\let\freeResponse\relax
\let\endfreeResponse\relax

%% \newtheoremstyle{ResponseStyle}{\topsep}{\topsep}%%% space between body and thm
%% 		{\wedn\bfseries}                      %%% Thm body font
%% 		{}                              %%% Indent amount (empty = no indent)
%% 		{\wedn\bfseries}            %%% Thm head font
%% 		{}                              %%% Punctuation after thm head
%% 		{3ex}                           %%% Space after thm head
%% 		{\underline{\underline{\thmname{#1}}}}%%% Thm head spec
%% \theoremstyle{ResponseStyle}

\usepackage[tikz]{mdframed}
\mdfdefinestyle{ResponseStyle}{leftmargin=1cm,linecolor=black,roundcorner=5pt,
, font=\bsifamily,}%font=\wedn\bfseries\upshape,}


\ifhandout
\NewEnviron{freeResponse}{}
\else
%\newtheorem{freeResponse}{Response:}
\newenvironment{freeResponse}{\begin{mdframed}[style=ResponseStyle]}{\end{mdframed}}
\fi



%% attempting to automate outcomes.

%% \newwrite\outcomefile
%%   \immediate\openout\outcomefile=\jobname.oc
%% \renewcommand{\outcome}[1]{\edef\theoutcomes{\theoutcomes #1~}%
%% \immediate\write\outcomefile{\unexpanded{\outcome}{#1}}}

%% \newcommand{\outcomelist}{\begin{itemize}\theoutcomes\end{itemize}}

%% \NewEnviron{listOutcomes}{\small\sffamily
%% After answering the following questions, students should be able to:
%% \begin{itemize}
%% \BODY
%% \end{itemize}
%% }
\usepackage[tikz]{mdframed}
\mdfdefinestyle{OutcomeStyle}{leftmargin=2cm,rightmargin=2cm,linecolor=black,roundcorner=5pt,
, font=\small\sffamily,}%font=\wedn\bfseries\upshape,}
\newenvironment{listOutcomes}{\begin{mdframed}[style=OutcomeStyle]After answering the following questions, students should be able to:\begin{itemize}}{\end{itemize}\end{mdframed}}



%% my commands

\newcommand{\snap}{{\bfseries\itshape\textsf{Snap!}}}
\newcommand{\flavor}{\link[\snap]{https://snap.berkeley.edu/}}
\newcommand{\mooculus}{\textsf{\textbf{MOOC}\textnormal{\textsf{ULUS}}}}


\usepackage{tkz-euclide}
\tikzstyle geometryDiagrams=[rounded corners=.5pt,ultra thick,color=black]
\colorlet{penColor}{black} % Color of a curve in a plot



\ifhandout\newcommand{\mynewpage}{\newpage}\else\newcommand{\mynewpage}{}\fi

\title{Supersize me}

\author{Bart Snapp}

\begin{document}
\begin{abstract}
  We will investigate the physical effects of scaling.
\end{abstract}
\maketitle


\begin{listOutcomes}
\item Explain the relation between linear scaling and area scaling.
\item Explain the relation between linear scaling and volume scaling.
\item Explain the relation between linear scaling and weight scaling.
\item Apply understanding of scaling to evaluate thought experiments.
\end{listOutcomes}

This is Dr.\ Zoidberg:

\begin{center}
  \includegraphics[width=.3\textwidth]{zoidberg.png}
\end{center}

In the \textit{Futurama} episode \link[\textit{Anthology of Interest
    I}]{https://youtu.be/-IOXv1ysae4} Zoidberg is enlarged to be
approximately $100$ times his normal height. Let's use what we know
about scaling to think about this.

\mynewpage


 
\begin{question}
  Compute the following:
    \begin{enumerate}
    \item $\frac{\text{Dr.\ Zoidberg's enlarged
        height}}{\text{Dr.\ Zoidberg's normal height}}$
    \item $\frac{\text{Dr.\ Zoidberg's enlarged
        surface area}}{\text{Dr.\ Zoidberg's normal surface area}}$
    \item $\frac{\text{Dr.\ Zoidberg's enlarged
        volume}}{\text{Dr.\ Zoidberg's normal volume}}$
    \end{enumerate}
    In each case, EXPLAIN your reasoning.
    \begin{freeResponse}
      \begin{enumerate}
      \item Well,
        \[
        \frac{\text{Dr.\ Zoidberg's enlarged
            height}}{\text{Dr.\ Zoidberg's normal height}} = 100.
        \]
      \item And,
        \[
        \frac{\text{Dr.\ Zoidberg's enlarged
            surface area}}{\text{Dr.\ Zoidberg's normal surface area}} = 100^2 = 10000
        \]
        because when you scale linear lengths by a factor, you scale
        the area by the square of the factor.
      \item Finally,
        \[
        \frac{\text{Dr.\ Zoidberg's enlarged
            volume}}{\text{Dr.\ Zoidberg's normal volume}} = 100^3 = 1000000
        \]
        because when you scale linear lengths by a factor, you scale
        the volume by the cube of the factor.
      \end{enumerate}
    \end{freeResponse}
\end{question}
\mynewpage


\begin{question}
  Now estimate the following:
  \begin{enumerate}
    \item $\frac{\text{Dr.\ Zoidberg's enlarged
        weight}}{\text{Dr.\ Zoidberg's normal weight}}$
    \item $\frac{\text{Amount of Dr.\ Zoidberg's enlarged ``blood''}}{\text{Amount of Dr.\ Zoidberg's normal ``blood''}}$
    \item $\frac{\text{Weight of Dr.\ Zoidberg's enlarged
        shell}}{\text{Weight of Dr.\ Zoidberg's normal shell}}$
  \end{enumerate}
  In each case, DISCUSS how you arrived at your answer.
  \begin{freeResponse}
    \begin{enumerate}
    \item So,
      \[
      \frac{\text{Dr.\ Zoidberg's enlarged
        weight}}{\text{Dr.\ Zoidberg's normal weight}} = 100^3
      \]
      because volume determines weight!
    \item In this case:
      \[
      \frac{\text{Amount of Dr.\ Zoidberg's enlarged ``blood''}}{\text{Amount of Dr.\ Zoidberg's normal ``blood''}} = 100^3
      \]
      because volume determines the amount of blood!
    \item This one is tricky. If Zoidberg's shell is very thin, the answer will be closer to
      \[
      \frac{\text{Weight of Dr.\ Zoidberg's enlarged
          shell}}{\text{Weight of Dr.\ Zoidberg's normal shell}} = 100^2
      \]
      because the weight of a thin shell will be determined by its
      area. On the other hand, if there is a significant thickness to the shell then  the answer will be closer to:
      \[
      \frac{\text{Weight of Dr.\ Zoidberg's enlarged
          shell}}{\text{Weight of Dr.\ Zoidberg's normal shell}} = 100^3
      \]
      because volume determines weight.
    \end{enumerate}
  \end{freeResponse}
\end{question}
\mynewpage



\begin{question}
  I think if I were made $100$ times larger:
  \begin{enumerate}
  \item I would suffocate!
  \item Also I would not be able to lift my arms or legs.
  \end{enumerate}
  Can you EXPLAIN why I might think that?

  \begin{freeResponse}
    \begin{enumerate}
    \item The amount of blood one has is determined by volume, but the
      amount of oxygen our lungs can absorb is determined by surface
      area. So I would have $1000000$ times more blood, but only
      $10000$ times as much surface area to acquire oxygen. Hence it
      would be as if I was getting $1/100$ of the amount of air I
      needed.
    \item So, your muscles work by contracting ``strings,'' and
      stronger muscles have more of these strings. Hence the strength
      of a muscle is determined by its cross-sectional area. Since the
      weight of me, $100$ times larger is $1000000$ more, and my
      strength is only $10000$ more, I would be $100$ times weaker
      than I am. I don't think I would be able to move. 
    \end{enumerate}
  \end{freeResponse}
  
\end{question}



\end{document}
